% This is the ADASS_template.tex LaTeX file, 24 Oct 2025.
% It is based on the ASP general author template file, but modified to reflect the specific
% requirements of the ADASS proceedings.

% This is the aspauthor.tex LaTeX file
% Copyright 2021, Astronomical Society of the Pacific Conference Series
% Revision:  29 November 2021

% Content authored by M. Buttfield-Addison
% Revision:  08 November 2026
\documentclass[11pt,twoside]{article}

\usepackage{asp2023}

\aspSuppressVolSlug
\resetcounters

\bibliographystyle{asp2021}

\markboth{Buttfield-Addison}{Interrogating Processes in Australian Radio Astronomy Data Production}

\begin{document}

\title{Interrogating Processes in Australian Radio Astronomy Data Production}

\author{M. Buttfield-Addison}
\affil{University of Tasmania, Hobart, Tasmania, Australia \email{m.buttfield-addison@utas.edu.au}}

\paperauthor{Mars Buttfield-Addison}{m.buttfield-addison@utas.edu.au}{0000-0001-8512-2434}{University of Tasmania}{School of Information and Communication Technology}{Hobart}{Tasmania}{7000}{Australia}

%\aindex{Buttfield-Addison,~M.}

\begin{abstract}
Australian radio astronomy is entering an era of unprecedented scale. To plan for the future, it is critical to understand the socio-technical factors that affect radio data in Australia today, how they may change, and how this might affect data outputs.

Targeted interviewing compared the experiences of individuals working in software development, data analysis, operations, and leadership across organisations around Australia. In discussing a range of topics surrounding their work, participants raised valid and consistent concerns for the future that demand consideration.
\end{abstract}

% == Subject Index
%\ssindex{conference!ADASS 2026}
%\ssindex{organisations!UTAS}
% TODO double-check these with index list before final

% == ASCL Index
%\ooindex{FOOBAR, ascl:1101.010}
% TODO replace and double-check this before final

\section{Justification}

Radio science in Australia is at a critical juncture. The SKA-Low is nearing completion, and other instruments around the country have seen significant capability upgrades. There are now more components, procedures, and personnel involved in the production of radio data products than ever before.

As we near this new era of exa-scale astronomy, systems are not the only things that need to adapt; methods of governance, information-sharing, and standardisation must also be considered. But we can’t plan how these need to change if we don’t well understand how they work now. Of course, clever people from institutions all over the world have collaborated to plan for these things---but available information about current operations that such plans are based upon tends to skew towards historical, international, and late-career perspectives. This limited view of operations may fail to capture Australia's unique context within the field of radio astronomy operations.

Herein, I seek to address this by interrogating the complex interrelations of technology, information, and people that are responsible for capturing, processing, and distributing radio data in Australia today, to ensure that next-generation operations such as SKA-Low can be based on current and contextual information, and set for success.

\subsection{Prior Research}

As a nation, Australia is blessed with several prolific historians of astronomy, and this makes it easy to find \textit{historical} records of astronomy operations and culture. But when it comes to contemporary records, a conspicuous gap appears---particularly regarding data \textit{producers} rather than just the data \textit{consumers} who rely on them. Prior research divides almost uniformly into work that was Australia-specific and partially first-hand but not current \citep[e.g.][]{orchiston2021golden,schilizzi2024square}, and work that surveyed the field in the right way but was not Australian and focused on data consumers rather than producers \citep[e.g.][]{ZUIDERWIJK2019228,Borgman2021From}.

Separately, a glut of documentation exists on organisational planning done by various parties at various points. But these are just that---plans---and there is no evidence to suggest these plans were adhered to or kept updated. Earlier work suggested they were unlikely to be, and results herein confirm that was a sound assumption.

\subsection{The Anecdotal Piece}

As part of my introduction to the field of radio astronomy instrumentation, I worked with scientists on various pieces of software to process data. In doing so, I was regularly allocated tasks with caveats; ``to use data from \texttt{[observatory]}, you need to \texttt{[manual re-interpretation]} the data'' was said in different ways about many a data source.

It became clear that even where interoperability protocols were agreed upon, their implementation could vary wildly. But the issues this should cause were being covered over in intermediate stages by specific individuals applying insider knowledge as had been communicated to them via informal social networks. So no-one downstream was any the wiser that non-compliance was occurring, nor that the correctness and compatibility of their data now relied on key individuals' information and ongoing involvement.

I knew that effective discovery and understanding of such technical challenges---if they did exist elsewhere in the field, as I predicted---would too require specialised knowledge and contacts to succeed. Having built some of these, I was in a unique position to interrogate the issue.

I began first by conducting an anonymous survey of VLBI operators around the world about their adoption and validation of a raw data format \citep[see][]{buttfield2022vdif}, and later a document review focused on advertised versus verified capabilities of various instruments (in prep.), but these only reinforced two main things: first, that issues were broader than such focused inquiries could uncover, and second, that common issues arising in Australia were more pronounced than those in more populous regions.

\section{Methods}

This project forms part of a larger body of work that uses a mixed-methods approach to examine the same organisations, teams, instruments, and phenomena from different angles to compare and control for (as much as is possible) participant and researcher bias. However, the focus of this paper is on the current ongoing phase: expert interviews.

The first task was to establish specific individuals to speak to, and what to speak to them about, that could provide the informative, accurate, and relevant insight into the operations we seek to understand. Beginning with a list of organisations and roles that were doing such work in the field, I looked to public information online---such as staff listings---to identify individuals that held relevant titles within each. 

\textbf{Organisations:} CSIRO, University of Tasmania, CIRA, SKAO, AusSRC, ADACS, ICRAR, Pawsey, and others. \textbf{Roles:} software or firmware development, operations, data analysis or curation, data user support, leadership of any of the above.

The first iteration of the list had \textasciitilde45 members. As interviews proceeded, participants were asked to suggest peers they thought relevant; today, there are almost 100.

After 12 informal sessions in the background phase, formal interviews began. At this early juncture, the collected data comprises 18 individual interviews and 1 site visit.

The initial list of topic areas was compiled primarily from the related literature identified earlier, such that results could be compared to this previous work. Focus areas spanned technical development, collaboration, prioritisation, validation, upskilling, and knowledge transfer. Background sessions then identified topical factors, including broader issues like the impacts of generative AI or COVID-era workplace closures.

Analysis of the interview transcripts was done with a combination of standard qualitative human research methods including open coding, member checking, and batch interviewing, with a view to evolve shared themes into expert personas, but this is beyond the scope of this paper to discuss in full.

\section{Early Findings}

Though significant in contribution, Australia’s radio astronomy industry has historically been small in size---driven by individuals with broad skills, deep institutional knowledge, and long tenures. Even as significant efforts like the SKA-Low increase, these relative qualities remain consistent in comparison to operations elsewhere in the world. 

Despite the experiences of interviewees in this small initial cohort differing significantly, the emerging themes of discussion (\textbf{in bold}) were notably consistent...

\textbf{Institutional Memory} was the strongest theme, woven through all other discussion. This is the argument that only through experience with specific processes and instruments does one build a familiarity and expertise, and this is then relied upon thereafter even if their team is unaware. This caused strong concern about \textbf{Succession}, as the pipeline for regular new personnel and the support for hands-on mentoring that is required to transfer skills and ensure continuity of expertise within the organisation had long ago fallen behind. Thus, in every organisation---including new ones like SKAO, because they had been so long in planning---there were individuals involved in critical processes despite either having retired or being well overdue to. The vast majority viewed their operations as rife with \textbf{Individual Points of Failure} in the form of uniquely connected and knowledgeable personnel. And often this manifested as \textbf{Issues Known Only to Some} and \textbf{Nothing is Written Down}.

\textbf{The Astronomer/Developer Divide} is a long-running phenomenon in astronomy where personnel are typically trained primarily as domain experts in astronomy theory or in software development practice, but not both. Consequently, teams must manage individuals from each background \textbf{Augmenting Skills on the Job} to be able to perform both kinds of work. In older teams, astronomers with self-taught coding skills dominated, but newer teams showed the opposite. The \textbf{Inter-Personnel Skills Interface} was therefore one of the main challenges to contend with, and multiple teams believed the interface was becoming less formal and effective over time.

\textbf{Politics in Operations} refers to the influence of operations context on what is possible or encouraged in the technical domain. Discussion of what people were \textit{willing to fund} or \textit{promises never intended to be kept} dominated when participants were asked why they approached particular projects the way they did. Promised timelines that were never possible, physical infrastructure of compromised quality due to the funding landscape, and constant mention of the support to build but never to operate or maintain. This then led to \textbf{Hidden Networks} that were used to transfer everything from hardware and software to knowledge between organisations and even countries where formal connections were inefficient or impossible. This ranged from ``ex-colleagues now work at different observatories and send expensive parts back and forth'' to ``instrument B is compatible with instrument A because the same guy worked at both''. This worked to great benefit in the short term but placed strong \textbf{Limitations to Scale}.

\textbf{Responsive, Not Proactive, Change} was evident in both infrastructure and processes. Over time, systems that had grown rather than been designed led to critical components that were ill-suited and poorly understood. Worse, the work required to inevitably overhaul grew over time. Participants reported reliance on funding and collaboration models that had failed to keep track with infrastructure scale and necessities.

At the same time, \textbf{The Service Model} has completely changed the nature of astronomy compared to an era where the same teams both collected and analysed data. Data producers state that, in many cases, \textbf{Users Don't Want To Know} or deal with the nuances of what came before the data repository. So if they are going to assume the data is source-agnostic, then for repeatable and accurate science to get done, it must be.

\section{Closing}

As more interviews are conducted, a critical mass of perspectives will allow more in-depth discussion without compromising anonymisation. But even these early themes paint a clear picture of passionate people working under incredible pressures beyond the technical complexity of their domain.

Of course, to the reader demographic of this paper, none of this is surprising. But value lies in identifying and communicating challenges in a way that those who fund or oversee our operations can understand. And ideally, for that to inform organisation and process changes that might ease the  burden of these experts for whom overwork, individual responsibility, and invisible labour is unfortunately commonplace.

\subsection{At ADASS XXXVI}

This work was presented at ADASS XXXVI in Perth, Australia in November 2026. This was beneficial to improve awareness and recruitment for the project, however, publicising of early results was primarily motivated by the potential to be part of ongoing discussions; many a delegate reinforced that several factors (beyond scale of operations) were forcing change---whether it be generative AI models in software development and knowledge discovery, barriers to hardware acquisition or data sharing across borders, or the emerging focus on multi-spectral/multi-repository data fusion.

\subsection{Invitation to Contribute}

Recruitment for this interviewing process will be ongoing until at least mid-2027. If you are interested, see \href{https://themartianlife.com/research-interview}{\texttt{themartianlife.com/research-interview}}. And in the spirit of open dialogue, I happily invite comments, questions, or criticism from any member of the community at \href{mailto:mail@themartianlife.com}{\texttt{mail@themartianlife.com}}.

\bibliography{C603}

% if we have space left, we might add a conference photograph here. Leave commented for now.
% \bookpartphoto[width=1.0\textwidth]{foobar.eps}{FooBar Photo (Photo: Any Photographer)}

\end{document}
